% Part: first-order-logic
% Chapter: model-theory
% Section: theory-of-m

\documentclass[../../../include/open-logic-section]{subfiles}

\begin{document}

\olfileid{mod}{bas}{thm}

\section{结构的理论}

每个结构~$\Struct{M}$ 都使一些语句为真，另一些语句为假。所有在其中为真的语句组成的集合，称为它的\emph{理论}。这个集合确实是一个（演绎封闭的）理论，因为由它所语义后承的任何语句，必然在它的所有模型（包括~$\Struct{M}$）中为真。

\begin{defn}
  给定结构~$\Struct M$，$\Struct{M}$ 的\emph{理论}是指在 $\Struct{M}$ 中为真的所有语句构成的集合 $\Theory{M}$，即 $\Theory{M} =
  \Setabs{!A}{\Sat{M}{!A}}$。
\end{defn}

我们也非正式地使用“理论”一词，来指代那些具有预期解释的语句集合，无论其是否演绎封闭。

\begin{prop}
对任意 $\Struct{M}$，$\Theory{M}$ 都是完备的。
\end{prop}

\begin{proof}
对任意语句~$!A$，要么 $\Sat{M}{!A}$，要么 $\Sat{M}{\lnot !A}$，因此要么 $!A \in \Theory{M}$，要么 $\lnot !A \in \Theory{M}$。
\end{proof}

\begin{prop}\ollabel{prop:equiv}
  如果对每个 $!A \in \Theory{M}$，都有 $\Struct{N} \models !A$，那么 $\Struct{M} \elemequiv \Struct{N}$。
\end{prop}

\begin{proof}
由于对所有 $!A \in \Theory{M}$，都有 $\Sat{N}{!A}$，故 $\Theory{M} \subseteq
\Theory{N}$。对任意语句 $!A$，若 $\Sat{N}{!A}$，则 $\Sat/{N}{\lnot !A}$，于是 $\lnot !A
\notin \Theory{M}$。由于 $\Theory{M}$ 是完备的，故 $!A \in
\Theory{M}$。因而 $\Theory{N} \subseteq \Theory{M}$，由此可得 $\Struct{M} \elemequiv \Struct{N}$。
\end{proof}

\begin{rem}\ollabel{remark:R}
  考虑结构 $\Struct{R} = \langle\Real, <\rangle$，其论域为实数集 $\Real$，其语言仅包含一个被解释为实数上 $<$ 关系的二元谓词符号。显然，$\Struct{R}$ 是不可数的；然而，由于 $\Theory{R}$ 显然是一致的，根据 Löwenheim--Skolem 定理，它有一个可数模型，记作 $\Struct{S}$；并且根据 \olref{prop:equiv}，有 $\Struct{R}
  \equiv \Struct{S}$。此外，由于 $\Struct{R}$ 和 $\Struct{S}$ 并不同构，这表明 \olref[iso]{thm:isom} 的逆命题在一般情况下不成立。
\end{rem}

\end{document}